% SEGMENT OLP-0700-S01
% Part: proof-theory
% Chapter: sequent-calculus
% Section: introduction

\documentclass[../../../include/open-logic-section]{subfiles}

\begin{document}

\olfileid{pt}{seq}{int}

\olsection{அறிமுகம்}

தொடரணிக் கணியங்கள், 1930-களில்
கெர்ஹார்ட் ஜென்ட்சன் அறிமுகப்படுத்திய
!!{proof} முறைமைகள். கூற்றுகளை
!!{proving} செய்வதற்கான
நடைமுறை வழியாக இருப்பதே
அவற்றின் முதன்மை நோக்கமல்ல.
மாறாக, !!{proof}s கொள்ளக்கூடிய
கட்டமைப்புகள் மற்றும் அவற்றின்
பண்புகளைப் பற்றிய சில முடிவுகளை
நிறுவ ஜென்ட்சன் அவற்றைப்
பயன்படுத்தினார். இப்படிப்பட்ட
கேள்விகளையே நிறுவல் கோட்பாடு
ஆராய்கிறது.

% SEGMENT OLP-0700-S02
பெயர் குறிப்பதுபோல்,
தொடரணிக் கணியங்கள்
!!{formula}s மீது அல்ல,
\emph{தொடரணிகள்} மீது
செயல்படுகின்றன. ஒரு தொடரணி,
!!{formula}s இன்
இரு வரிசைகள் அல்லது
இரு பன்மைக் கணங்களால்
ஆனது. ஜென்ட்சனின்
மூல முறைமைகளான
\Log{LK}, \Log{LJ}
வரிசைகளைப் பயன்படுத்தின;
இன்றைய நிறுவல்
கோட்பாட்டாளர்கள்
பன்மைக் கணங்களைப்
பயன்படுத்த விரும்புகிறார்கள்.
ஒரு பன்மைக் கணம்
சாதாரணக் கணத்தைப் போன்றது;
ஆனால் அதில் ஓர்
!!{element} ஒன்றுக்கு
மேற்பட்ட முறை
இடம்பெறலாம். இருப்பினும்,
சாதாரணக் கணத்தைப் போலவே,
!!{element}s இன்
வரிசை அதில் பொருட்டல்ல.
எனவே $\{!A, !A, !B\}$
மற்றும் $\{!A, !B, !A\}$
ஒரே பன்மைக் கணத்தைக்
குறிக்கின்றன. ஆனால்
அது $\{!A, !B\}$
என்பதிலிருந்தும்
$\{!A, !A, !B, !B\}$
என்பதிலிருந்தும் வேறுபடும்:
$!A$, $!B$ ஆகியவற்றின்
நிகழ்வெண்கள் அவற்றில்
வேறுபடுகின்றன.

% SEGMENT OLP-0700-S03
நிறுவல் கோட்பாட்டில்
!!{formula}s இன்
பன்மைக் கணங்களை
வழக்கமாகப் பெரிய
கிரேக்க எழுத்துகளால்
குறிக்கிறார்கள். ஆகவே
$\Gamma$, $\Delta$,
$\Pi$, $\Lambda$,
அல்லது $\Theta$
என்று எழுதும்போது,
எந்தத் தருக்கத்தை
ஆராய்கிறோமோ அதன்
!!{formula}s இன்
பன்மைக் கணங்களையே
குறிக்கிறோம். மேலும்
வழக்கமாக,
$\tuple{\Gamma, \Delta}$
என்று எழுதுவதற்குப்
பதிலாக இரு கூறுகளையும்
ஒரு தொடரணிக்
குறியால் பிரிக்கிறார்கள்;
இங்கு~$\Sequent$
ஐப் பயன்படுத்துகிறோம்.

\begin{defn}[தொடரணி]
\emph{தொடரணி} என்பது
\[
\Gamma \Sequent \Delta
\]
என்ற வடிவிலான கோவை;
இங்கு $\Gamma$, $\Delta$
இரண்டும் முடிவுடைய,
வெறுமையாகவும் இருக்கக்கூடிய
!!{formula}s இன்
பன்மைக் கணங்கள்.
$\Gamma$ \emph{முன்தொடர்}
எனவும், $\Delta$
\emph{பின்தொடர்} எனவும்
அழைக்கப்படுகின்றன.
\end{defn}

% SEGMENT OLP-0700-S04
$!G$ என்பது~$\Gamma$ வில்
உள்ள !!{formula}s இன்
இணையல் எனவும்
(வெற்று இணையலை~$\ltrue$
எனக் கொண்டு),
$!D$ என்பது~$\Delta$
வில் உள்ள !!{formula}s
இன் பிரிப்பிணைவு எனவும்
(வெற்றுப் பிரிப்பிணைவை
~$\lfalse$ எனக் கொண்டு)
இருக்கட்டும். அப்போது
தொடரணி $\Gamma \Sequent
\Delta$, ஏதேனும்
கட்டமைப்பு~$\Struct{M}$
இல் மெய்யாக இருப்பதும்
$!G \lif !D$
மெய்யாக இருப்பதும்
ஒரே நிபந்தனை.
பாரம்பரியத் தருக்கத்தில்
இதன் பொருள்:
$\Gamma$ வில் உள்ள
!!{formula}s இல்
குறைந்தது ஒன்று
பொய்யாகவோ,
$\Delta$ வில் உள்ள
!!{formula}s இல்
குறைந்தது ஒன்று
மெய்யாகவோ இருக்கும்.

% SEGMENT OLP-0700-S05
$\Gamma$ என்பது
!!{formula}s இன்
பன்மைக் கணமெனில்,
$\Gamma$ வில் $!A$ ஐச்
சேர்ப்பதன் விளைவை
$\Gamma, !A$
(அல்லது $!A, \Gamma$)
என்று எழுதுகிறோம்.
$\Delta$ உம்
!!{formula}s இன்
பன்மைக் கணமெனில்,
$\Gamma, \Delta$
என்பது அவ்விரு
பன்மைக் கணங்களின்
ஒன்றியம்: $\Gamma$
வில் $!A$ இன்
\emph{நிகழ்வெண்}~n
(அதாவது, $!A$
இன் $n$ பிரதிகள்)
என்றும், $\Delta$
வில் அதன்
நிகழ்வெண்~$m$
என்றும் இருந்தால்,
$\Gamma, \Delta$
வில் $!A$ இன்
நிகழ்வெண்~$n+m$.
ஒரு தொடரணியின்
ஒரு பக்கத்திலுள்ள
பன்மைக் கணம்
வெறுமையாக இருந்தால்,
வழக்கமாக அந்த இடத்தை
வெறுமையாகவே விடுகிறோம்.
எடுத்துக்காட்டாக,
$\Sequent !A$
என்பது வெற்றான
முன்தொடரையும்,
நிகழ்வெண்~$1$
உடன் $!A$ வரும்
பின்தொடரையும்
கொண்ட தொடரணி.

% SEGMENT OLP-0700-S06
தொடரணிக் கணியத்தில்
!!^a{proof} என்பது
தொடரணிகளால் ஆன
மரம். அதன் உச்சியில்
உள்ள தொடக்கத்
தொடரணிகள்,
ஆராயப்படும் முறைமையின்
அடிக்கோள்களாக
இருக்க வேண்டும்.
ஒரு தொடரணி,
ஒன்று அல்லது இரண்டு
தொடரணிகளுக்குக் கீழே
இருந்தால், அந்த
முறைமையின் அனுமான
விதியால் அவற்றிலிருந்து
சரியாகப் பெறப்பட
வேண்டும். முதலில்,
பாரம்பரியத் தருக்கத்திற்கான
இரு வேறு தொடரணி
முறைமைகளான
$\Log{G1c}$,
~$\Log{G3c}$
ஆகியவற்றைப் பார்ப்போம்.
அவற்றின் அடிக்கோள்களும்
விதிகளும்
\cref{pt:seq:int:tab:G1c,pt:seq:int:tab:G3c}
இல் தரப்பட்டுள்ளன.
\oliflabeldef{fol:seq::chap}{ (ஜென்ட்சனின்
மூல முறைமை \Log{LK},
\olref[fol][seq][]{chap}
இல் விவரிக்கப்பட்டுள்ளது.)}{}

\subfile{rules-G1c}
\subfile{rules-G3c}

% SEGMENT OLP-0700-S07
இவ்விரு முறைமைகளுக்கு
இடையே நுட்பமான
வேறுபாடுகள் உள்ளன;
அவற்றால் நிறுவல்
கோட்பாட்டுச்
சிறப்பியல்புகளும்
வேறுபடுகின்றன.
\Log{G1c} இல்
$!A \Sequent !A$
என்ற அடிக்கோள்கள்
உள்ளன; வேறு
!!{formula}s
அவற்றில் இருக்க
அனுமதி இல்லை.
\Log{G3c} இல்
$!C, \Gamma \Sequent
\Delta, !C$
என்ற வடிவிலான
அடிக்கோள்கள் உண்டு;
$\Gamma$,
~$\Delta$ வில்
எந்தப் பக்க
!!{formula}s உம்
இருக்கலாம். ஆனால்
இரு பக்கங்களிலும்
வரும் !!{formula}~$!C$
\emph{அணுவாக}
இருக்க வேண்டும்.
\Log{G1c} இல்
$\LeftR{\land}$,
$\RightR{\lor}$
ஆகிய ஒவ்வொரு விதிக்கும்
இரண்டு வடிவங்கள்
உள்ளன; \Log{G3c}
இல் ஒவ்வொன்றுக்கும்
ஒரே வடிவம்தான்.
மேலும் \Log{G1c} இல்,
“சுருக்கல்”
(\LeftR{\Contraction},
\RightR{\Contraction})
மற்றும் “தளர்த்தல்”
(\LeftR{\Weakening},
\RightR{\Weakening})
என்ற கூடுதல்
“கட்டமைப்பு விதிகள்”
உள்ளன.

% SEGMENT OLP-0700-S08
\Log{G1c}, \Log{G3c}
இரண்டும் பாரம்பரியத்
தருக்கத்திற்குச்
சரியானவையும்
நிறைவானவையும் ஆகும்.
அதாவது, இம்முறைமைகளில்
!!{provable} ஆன
ஒவ்வொரு தொடரணியும்
ஒவ்வொரு !!{structure}
இலும் மெய்; ஒவ்வொரு
!!{structure} இலும்
மெய்யான ஒவ்வொரு
தொடரணிக்கும்
!!a{proof} உண்டு.
ஆகவே ஒரு தொடரணிக்கு
!!a{proof} \emph{உண்டா}
என்பதைத் தருக்கத்தின்
வேறு முறைகளாலும்
(எடுத்துக்காட்டாக,
மாதிரிக் கோட்பாட்டு
முறைகளாலும்)
பதிலளிக்கலாம்.
மேலும் இவ்விரு
முறைமைகளும் ஒவ்வொரு
!!{structure} இலும்
மெய்யான தொடரணிகளையே
துல்லியமாக !!{prove}
செய்வதால், அவை
ஒரே தொடரணிகளை
நிறுவுகின்றன.

% SEGMENT OLP-0700-S09
ஆனால் நிறுவல்
கோட்பாடு, ஒன்று
\emph{நிறுவத்தக்கதா}
என்பதில் மட்டுமன்றி,
\emph{எப்படி}
நிறுவப்படுகிறது
என்பதிலும் அக்கறை
கொண்டுள்ளது.
“ஒரு குறிப்பிட்ட
விதியை எப்போதும்
தவிர்க்க முடியுமா?”,
“புதிய தொடரணிகள்
நிறுவத்தக்கதாகாமல்
இந்த விதியை
முறைமையில் சேர்க்க
முடியுமா?”,
“ஒரு முறைமையின்
!!{proof}s ஐ
மற்றொரு முறைமையின்
!!{proof}s ஆக
மாற்ற முடியுமா?”
போன்ற கேள்விகளை
நிறுவல் கோட்பாட்டாளர்கள்
ஆராய்கிறார்கள்.
இத்தகைய கேள்விக்கு
“ஆம்” என்ற விடையை
நிறுவும்போது,
!!{proof}s இன்
சிக்கல்தன்மையின்
மீதோ, அதற்கு
இணையான அவற்றின்
கட்டமைப்பின்
மீதோ தொகுத்தறிதல்
செய்வது வழக்கம்.
இதனால்,
\Log{G1c} ஒரு
தொடரணியை நிறுவினால்
\Log{G3c} உம்
அதே தொடரணியை
நிறுவும் என்பது
போன்ற உண்மைக்கான
நிறுவல் மட்டுமல்ல,
\Log{G1c} இன்
!!{proof}s ஐ
\Log{G3c} இன்
!!{proof}s ஆக
மாற்றும் செயன்முறையும்
கிடைக்கும்.

% SEGMENT OLP-0700-S10
நிறுவல் கோட்பாட்டின்
மைய முடிவுகளில்
ஒன்று \emph{வெட்டு
நீக்கத் தேற்றம்}.
வெட்டு விதியை,
\[
\Axiom$\Gamma \fCenter \Delta, !A$
\Axiom$!A, \Pi \fCenter \Lambda$
\RightLabel{\Cut}
\BinaryInf$\Gamma, \Pi \fCenter \Delta, \Lambda$
\DisplayProof
\]
தொடரணி முறைமைகளுடன்
சேர்த்தாலும் எந்தத்
தொடரணிகள்
!!{provable}
என்பது மாறாது
என்று இத்தேற்றம்
காட்டுகிறது.
மேலும்,
\Log{G1c}
(அல்லது \Log{G3c})
விதிகளோடு \Cut{}
ஐயும் பயன்படுத்தும்
எந்த !!{proof}
ஐயும், \Log{G1c}
(அல்லது \Log{G3c})
விதிகளை மட்டும்
பயன்படுத்தும்
!!{proof} ஆக
மாற்றும் முறையை
அதன் நிறுவல்
தருகிறது. வேறு
சொற்களில், இவ்வாறு
பயன்படுத்திய
!!{proof}s இலிருந்து
\Cut{} விதியை
இந்தச் செயன்முறை
“நீக்குகிறது”;
இதனால்தான்
அப்பெயர்.

\end{document}
